% 来源：OEBPS/Text/01_37.xhtml（OPF spine 顺序；仅第一册）
\cleardoublepage
\BookEntry[本杰明·格雷厄姆《聪明的投资者》]{本杰明·格雷厄姆《聪明的投资者》}

\BookCover{images/书单22聪明的投资者.jpg}{[美] 本杰明·格雷厄姆著| 王中华、黄一义译| 人民邮电出版社2016年版}

\begin{BookSummaryBox}
{\centering \textbf{“现代证券分析之父”、价值投资理论奠基人\\ 格雷厄姆撰写的股市“圣经”}\par}
\end{BookSummaryBox}

\subsubsection*{这本书谈什么？}

《聪明的投资者》（第4版）在保留格雷厄姆1973年原著的基础上，由贾森·兹威格在近40年特别是世纪之交全球股市大动荡的现实之后，进一步检验和佐证了价值投资理论。
\par

这本书主要面向个人投资者，旨在对普通人在投资策略的选择和执行方面提供一些帮助和指导。首先，该书明确了“投资”与“投机”的区别，指出聪明的投资者当如何确定预期收益；其次，以介绍防御型投资者与积极型投资者的投资组合为策略，阐释了投资者如何应对市场波动；随后，本书还对基金投资、投资者与投资顾问的关系、普通投资者证券分析的一般方法、防御型投资者与积极型投资者的证券选择等问题进行了详细说明；最后，作者通过分析几组具体案例，论述了股息政策和作为投资中心思想的“安全性”问题。
\par

这不是一本教人“如何成为百万富翁”的书，但主要目标是讲述规范的投资原理和投资态度，以防投资者陷入一些经常性的错误中。
\par

\subsubsection*{芒格为什么谈起这本书？}

对芒格投资观念产生重要影响的是两位“本杰明”先生，一位是富兰克林，另一位就是格雷厄姆。格雷厄姆的《聪明的投资者》中最具有生命力的观念之一就是“市场先生”。在这位享誉华尔街的投资思想家看来，一般情况下，市场先生是一个脾气温和、头脑理智的家伙，但有时候也会受到非理性的恐惧或贪婪的驱使。
\par

因此，格雷厄姆提醒投资者，对于股票的价值，要亲自作出客观的判断。同样，在芒格看来，即使是最有能力、最有干劲的人，也不总是能基于理性而作出决定。
\par

\subsubsection*{《聪明的投资者》作者本杰明·格雷厄姆如是说}

“聪明的定义是：要有耐心，要有约束，并渴望学习；此外，还必须能够驾驭你的情绪，并能够进行自我反思。”
\par

“如果投资者自己因为所持证券市场价格不合理的下跌而盲目跟风或过度担忧的话，那么，他就是不可思议地把自己的基本优势转变成了基本劣势。投资活动并非要在别人的游戏中打败他们，而是要在自己的游戏中控制好自己。”
\par

“所有美好的东西,都是既罕见又复杂。即使是聪明的投资者，也可能需要很强的意志力来防止自己的从众行为。”
\par
